\begin{table}[H]
\centering
\caption{LP: residual calculation, sector 10}
\GOneLabel{tab:deep-lp-worked}
\begingroup\fontsize{9.2}{11.2}\selectfont
\setlength{\tabcolsep}{4.3pt}\renewcommand{\arraystretch}{1.1}
\begin{tabular}{lrr}
\toprule
Quantity & 2011 & 2022 \\
\midrule
Predicted records & 318 & 343 \\
Mean $y$ & 13.7219 & 14.9531 \\
Mean $l$ & 3.8132 & 3.6219 \\
Mean $k$ & 13.4471 & 14.8332 \\
$\widehat\beta_L$ & 0.5465 & 0.5465 \\
$\widehat\beta_K$ & 0.0867 & 0.0867 \\
$\widehat\beta_L\overline l$ & 2.0838 & 1.9792 \\
$\widehat\beta_K\overline k$ & 1.1658 & 1.2860 \\
Mean uncentered residual & 10.4723 & 11.6878 \\
Pooled fitted reference mean & 11.1596 & 11.1596 \\
Mean centered residual & -0.6873 & 0.5283 \\
Exponential of mean centered log & 0.5029 & 1.6960 \\
Mean exponentiated index & 2.1457 & 8.8368 \\
\bottomrule
\end{tabular}
\endgroup
\appnotes{We show the first and last prediction years with releasable observed group means in the first estimated sector that supplies such groups. Each year uses all eligible primary predictions in that group, including observations outside dynamic estimation. We first log individual inputs, then average the logs. Every year uses the same pooled sector coefficients and pooled dynamic-sample reference mean. The exponential of the mean centered log is the geometric mean of the index; the last row is its arithmetic mean. No individual company record supplies this example.}{G0.5 v3 modified manufacturing records with fictional company links; authors' calculations.}
\end{table}
